\Subhead{2.10.2}{Обращение волнового фронта}

Обычное зеркало меняет знак нормальной компоненты волнового вектора, но не обращает произвольную поперечную фазу. Фазосопрягающее зеркало создаёт встречное поле с комплексной амплитудой $A^*(x,y)$. Для $A=|A|\e^{\ii\phi(x,y)}$ это меняет знак аберрационной фазы и при обратном проходе через ту же стационарную взаимную среду может компенсировать её.

Один механизм — вырожденное четырёхволновое смешение: два встречных накачивающих поля и сигнал создают поляризацию
\begin{equation}
 P_c\propto\chi^{(3)}E_{p1}E_{p2}E_s^*,\quad
 \omega_c=\omega_{p1}+\omega_{p2}-\omega_s,\quad
 \bm k_c=\bm k_{p1}+\bm k_{p2}-\bm k_s.
\end{equation}
При равных частотах и встречных накачках $\bm k_c=-\bm k_s$. Другой механизм — вынужденное рассеяние назад, например на акустической волне; тогда сопряжённая волна обычно имеет частотный сдвиг. Конечные полоса, апертура, шум и изменение среды ограничивают компенсацию. Фазовое сопряжение монохроматического поля не тождественно полному обращению времени всех степеней свободы среды.
